% CLASS 3 -- the notes' section 1.3, sec.pointProcesses, COMPRESSED.
% Budget 75 min; planned 75 over 18 frames, so no headroom.
% Overrunning, cut F3.13 (fold def.relaxationTime into F3.12 as one line and drop its
% picture), then F3.3 (assert the survival function inline on F3.15 and F3.16).
% Do NOT cut F3.12: section 1.4.3 is unstateable without prop.forwardMean.
%
% DROPPED BY DESIGN: lemma.perronFrobenius; the whole second-order block (sec.secondOrder,
% thm.lyapunov, lemma.lyapunovUniqueness, eq.stationaryCovariance, eq.stateProductJump,
% eq.lyapunov, eq.explicitFormulaeOfStationaryCovariance, eq.secondMomentODE); the scalar
% case (eq.scalarVariance, remark.scalarVariance); sub- and supercriticality
% (eq.supercriticalGrowth, remark.martingaleNotBounded, remark.strictBaseline);
% prop.meanState; remark.noSubcriticality; corol.constantClusterSize as a statement;
% remark.hurwitzNeedsNonNegativity; the Poisson example and eq.NTE_counting_proc.
% Perron-Frobenius parts i and iv are used silently on F3.11 and F3.12, in a subordinate
% clause each, and get no statement.
%
% LABELS. A label is defined once across the five class files.

\subsection{Point processes and self-excitation}

\begin{frame}
  \frametitle{Outline}
% 3 min. One \nocite per line -- a multi-key \nocite spanning lines breaks bibtex.
\nocite{DVJ08int}
\nocite{Haw71spe}
\nocite{HO74clu}
\nocite{DZ13exa}
% The two features of the data that force the apparatus, taken from section 1.1's own
% paragraph: continuous physical time, and the effect of one type of event on the time
% distribution of others. Say what is kept and what is skipped.
\end{frame}

\begin{frame}
  \frametitle{Intensity and compensator}
% 5 min. One line on \multiCountingProc, the ground process \groundProc and the labels
% \event[n]; name eq.arrivalTimesOfGroundProcess and eq.eventLabels without displaying them.
% Then def.compensator with its four conditions, the absolutely continuous case, and the
% conditional expectation. Closing: predictability is the whole content -- \intensity(t) is a
% function of the STRICT past, which is why every intensity we write is left-continuous.
\end{frame}

\begin{frame}
  \frametitle{The conditional survival function}
% 4 min. prop.conditionalSurvival and eq.conditionalSurvival. The filtration is the INTERNAL
% history, and on the event that no arrival occurs the right-hand side is measurable, so it
% is a survival function we can invert. Both simulators are this proposition, one run
% backwards and one forwards.
\end{frame}

\begin{frame}
  \frametitle{Meyer's time change}
% 4 min. thm.meyer1971 and eq.time-changed_intertimes. The three hypotheses as a list --
% enumitem is now loaded, so the notes' own list syntax compiles. Closing: hypothesis i is
% what buys the JOINT independence; the univariate statement holds without it.
\end{frame}

\begin{frame}
  \frametitle{Hawkes processes}
% 4 min. def.hawkesprocess, eq.intensity_ordHawkes, eq.branchingMatrix. The half-open
% interval is what makes the intensity predictable.
\end{frame}

\begin{frame}
  \frametitle{The order of the indices}
% 4 min. remark.indexOrder. \alert{The row is excited and the column excites.}
% The column sum is the expected number of direct offspring of one event of type \eone.
% A matrix and its transpose share a spectral radius, so a transposed kernel passes every
% stability test and produces a path of the right total rate: the convention is the only
% defence against it.
\end{frame}

\begin{frame}
  \frametitle{The cluster representation}
% 5 min. prop.clusterRepresentation. PICTURE: a time axis with immigrant marks, each
% immigrant's offspring hanging below as a two-level tree with one whole cluster shaded, and
% the same points projected onto the axis as the path an observer sees. The picture IS the
% statement that the clusters are unobservable.
% Closing beat, ours rather than the notes': a cluster is a run with a common ancestor, and
% it is not a run of same-signed events.
\end{frame}

\begin{frame}
  \frametitle{The exponential kernel and its state}
% 4 min. def.exponentialHawkes, eq.exponentialKernel, eq.decayedCounts, eq.GammaIsAOverBeta,
% eq.intensityFromState, eq.stateDynamics. A SINGLE \decay shared by every pair. The strict
% inequality in \decayedCounts is what makes the intensity predictable. One clause for
% remark.countsAgainstRates: \branchingMatrix is dimensionless and \excitation is a rate.
\end{frame}

\begin{frame}
  \frametitle{The state is Markov}
% 5 min. thm.markovState and eq.generatorOfDecayedCount. The integrability condition in
% words, not symbols. Closing: the past enters through \numEventTypes numbers, so updating
% costs O(1) per event against O(n^2), and a decay per pair would need a state of dimension
% \numEventTypes squared -- exercise.eventDependentDecays as a clause, not a frame.
\end{frame}

\begin{frame}
  \frametitle{Stationarity and the mean intensity}
% 4 min. def.stationaryHawkes then thm.stationarity and eq.stationaryIntensity. Stationary
% iff \branchingRatio < 1. The one-line sketch: \stationaryIntensity = \baseIntensity +
% \branchingMatrix\stationaryIntensity, and \branchingRatio < 1 inverts.
% Introduce \totalRate and \bar\baseIntensity here.
\end{frame}

\begin{frame}
  \frametitle{Cluster size and the endogenous fraction}
% 5 min. def.endogenousFraction, prop.clusterSize, remark.clusterBudget, all reusing their
% labels. The cluster size COUNTS THE ANCESTOR. \bar\baseIntensity \bar C = \totalRate is the
% cheapest consistency check an implementation admits.
% \alert{\branchingRatio is neither the endogenous fraction nor 1 - 1/\bar C.}
% The two non-generic conditions in words: constant column sums, or \baseIntensity a right
% Perron vector. The shipped numbers 0.682 against 0.60 and 3.15 against 2.5 come from the
% package, not from the notes -- say so on the frame.
\end{frame}

\begin{frame}
  \frametitle{The forward mean and the Hurwitz matrix}
% 6 min. THE LOAD-BEARING FRAME OF THE CLASS. def.hurwitzMatrix, prop.hurwitzCriterion
% (\hurwitzMatrix = \excitation - \decay I is Hurwitz iff \branchingRatio < 1), and
% prop.forwardMean with both displays and \meanResponseIntegral(\forecastHorizon).
% Say aloud that eq.forwardMeanCounts is the equation class 4's forecast is a contraction of.
% One clause for remark.transientAgainstStationary.
\end{frame}

\begin{frame}
  \frametitle{The mean response and the relaxation time}
% 4 min. corol.meanResponse and def.relaxationTime. PICTURE: the mean response against t for
% two types, one decaying monotonically and one RISING BEFORE IT FALLS (a column sum of
% \branchingMatrix exceeding one), with the relaxation time marked as a tangent intercept and
% not as a half-life. The response integral ties back to the cluster size. First to cut.
\end{frame}

\begin{frame}
  \frametitle{Two dimensionless groups}
% 4 min. def.dimensionlessGroups and remark.dimensionlessGroups. \totalRate/\decay, events
% per excitation timescale, decides whether self-excitation is visible at the rate the
% process is sampled; \totalRate \ofiWindow decides whether a window sum is a statistic.
% \alert{At \totalRate/\decay well below one a study measures a process in which there is
% correctly nothing to find.}
\end{frame}

\begin{frame}
  \frametitle{Ogata's thinning algorithm}
% 5 min. algo.ogata as a bare algorithmic environment -- `algorithm' is a float and will not
% place inside a frame. The bound is free: \excitation >= 0 makes \totalIntensity
% non-increasing between events, so no separate majorant is constructed, and the dominating
% value is the POST-JUMP one. Then eq.update_of_exponential_intensity, O(1) against O(n^2).
% Closing beat as remark.prePostJump: the pre-jump value enters the intensity and the
% post-jump value is carried forward; confusing them shifts every intensity by one event and
% produces a path that is plausible and wrong.
\end{frame}

\begin{frame}
  \frametitle{The exact scheme --- 1 of 2}
% 4 min. eq.totalIntensityBetweenEvents and eq.dzSurvival: the common decay makes the excited
% part of \totalIntensity one exponential, so the survival function is a PRODUCT and the wait
% is a minimum of two independent clocks, one immigration and one defective excitation.
% thm.dassiosZhao with eq.dzDraws. One clause of remark.commonDecay: with a decay per pair
% the two-factor form fails and one clock per decay rate is needed.
\end{frame}

\begin{frame}
  \frametitle{The exact scheme --- 2 of 2}
% 4 min. algo.dassiosZhao as bare algorithmic. Nothing is approximate and nothing is
% rejected. remark.defectiveClock: the excited factor carries finite mass D/\decay, the
% expected number of direct offspring the state still owes, and with probability
% exp(-D/\decay) it expires and the next event is an immigrant -- clipping biases the
% immigrant share. remark.typeDraw: the wait sees \excitation only through its column sums,
% and the type draw is where the full matrix re-enters.
\end{frame}

\begin{frame}
  \frametitle{The residual test}
% 5 min. prop.compensatorClosedForm -- the compensator exactly, from the state and the
% counts, with no numerical integration; the count is read before T and the state at it
% (remark.readPreJump, one clause). Then the test: rescaled inter-arrivals are Exp(1),
% \cite{Oga88sta}. The three hypotheses, and which of them fails on a RECORDED path.
% Closing: remark.marginsOnly. \alert{A suite of \numEventTypes passing tests has established
% \numEventTypes margins.}
\end{frame}
